%%%%%%%%%%%%%%%%%
% CV tailored for PRO BTP - Chef de Projet Organisation H/F
%%%%%%%%%%%%%%%%%

\documentclass[8pt,a4paper,ragged2e,withhyper]{altacv}

\geometry{left=1.25cm,right=1.25cm,top=.5cm,bottom=1.cm,columnsep=1.2cm}

\usepackage{paracol}
\usepackage{fontawesome5}
\usepackage{hyperref}

\iftutex 
  \setmainfont{Roboto Slab}
  \setsansfont{Lato}
  \renewcommand{\familydefault}{\sfdefault}
\else
  \usepackage[rm]{roboto}
  \usepackage[defaultsans]{lato}
  \renewcommand{\familydefault}{\sfdefault}
\fi

\definecolor{SlateGrey}{HTML}{2E2E2E}
\definecolor{LightGrey}{HTML}{666666}
\definecolor{DarkPastelRed}{HTML}{8AC0D9}
\definecolor{PastelRed}{HTML}{00B0DA}
\definecolor{GoldenEarth}{HTML}{B4AD0C}

\colorlet{name}{black}
\colorlet{tagline}{PastelRed}
\colorlet{heading}{DarkPastelRed}
\colorlet{headingrule}{GoldenEarth}
\colorlet{subheading}{PastelRed}
\colorlet{accent}{PastelRed}
\colorlet{emphasis}{SlateGrey}
\colorlet{body}{LightGrey}

\definecolor{violet}{HTML}{5A2582}
\definecolor{rouge}{HTML}{F53B3B}
\definecolor{noir}{HTML}{00232C}
\definecolor{bleu}{HTML}{00B0DA}
\definecolor{rose}{HTML}{F831A0}
\definecolor{jaune}{HTML}{FFEC00}
\definecolor{vert}{HTML}{34AD6C}

\renewcommand{\namefont}{\Huge\rmfamily\bfseries}
\renewcommand{\personalinfofont}{\footnotesize}
\renewcommand{\cvsectionfont}{\LARGE\rmfamily\bfseries}
\renewcommand{\cvsubsectionfont}{\large\bfseries}
\renewcommand{\cvItemMarker}{{\small\textbullet}}
\renewcommand{\cvRatingMarker}{\faCircle}

\input{../../_shared/pubs-num.tex}
\addbibresource{../../_shared/sample.bib}

\begin{document}

\name{BOILEAU Guillaume}
\tagline{Chef de Projet Organisation | Transformation, Processus, Reporting, Performance \& SI}

\photoL{2.8cm}{../../_shared/moi}

\personalinfo{%
  \email{guillaume.boileaupro@gmail.com}
  \phone{+33 6 59 94 85 84}
  \location{Alpes-Maritimes, FRANCE}
  \linkedin{guillaume-boileau-phd-891b81265}
  \researchgate{Guillaume-Boileau-2}
  \orcid{0000-0002-3576-6968}
}

\makecvheader
\vspace{-0.3cm}

\AtBeginEnvironment{itemize}{\small}
\columnratio{0.64}

\begin{paracol}{2}

\cvsection{Experience}

\cvevent{\textbf{Software Engineer - Project Coordination, Process Follow-Up \& Operational Reporting}}{VEV Platform Services France}{2025 -- 2026}{Valbonne, France}

\textbf{Context:} Digital services and operational software for EV charging infrastructure, involving platform services, data flows, hardware-connected systems and production constraints.

\textbf{Objective:} Support reliable delivery, operational follow-up, issue resolution and coordination between technical stakeholders in an industrial digital environment.

\begin{itemize}
  \item \textbf{Digital project environment:} Contributed to the development and maintenance of web and backend services using TypeScript, Angular and Node.js.
  \item \textbf{Operational coordination:} Supported task follow-up, prioritisation and coordination with technical stakeholders in an Agile delivery context.
  \item \textbf{Process follow-up:} Monitored operational data flows, service behaviour and interface continuity to support quality and service reliability.
  \item \textbf{Issue analysis:} Investigated production issues involving software behaviour, data exchanges and hardware-connected systems.
  \item \textbf{Performance and quality:} Analysed data inconsistencies, abnormal behaviours and service anomalies affecting operational continuity.
  \item \textbf{Tools and reporting:} Used Zenhub and collaborative tracking logic to support backlog follow-up, issue tracking and technical coordination.
  \item \textbf{Operational interfaces:} Worked at the interface between software services, field equipment constraints and operational business needs.
\end{itemize}

\divider

\cvevent{\textbf{Senior R\&D Engineer - Transformation of Workflows, Technical Coordination \& Performance Analysis}}{CNES, Laboratoire Lagrange, Observatoire de la C\^ote d'Azur}{2023 -- 2025}{Nice, France}

\textbf{Context:} Engineering and coordination activities for the LISA ESA/NASA space mission, involving complex workflows, multiple contributors, technical deliverables, validation campaigns and international stakeholders.

\textbf{Objective:} Structure, coordinate and improve software/data workflows under strong traceability, quality, performance and documentation constraints.

\begin{itemize}
  \item \textbf{Project structuring:} Structured complex analysis and validation workflows for the \textbf{LISA ESA/NASA space mission}, involving several contributors and technical stakeholders.
  \item \textbf{Process improvement:} Designed and improved methods for model integration, output comparison, consistency checks and technical result consolidation.
  \item \textbf{Technical roadmap support:} Translated scientific and technical needs into development priorities, validation tasks and actionable work packages.
  \item \textbf{Performance analysis:} Analysed results, limitations, uncertainty sources and performance indicators to support interpretation and decision-making.
  \item \textbf{Reporting and governance support:} Produced structured technical notes, validation summaries, progress updates and material for technical reviews.
  \item \textbf{Coordination:} Coordinated contributors, research engineers, technicians and Master's thesis interns around shared technical objectives.
  \item \textbf{Documentation:} Used Confluence to support technical documentation, knowledge sharing, coordination notes and follow-up of validation activities.
  \item \textbf{Risk and issue identification:} Identified inconsistencies, deviations, weak points and operational limitations in complex software/data workflows.
  \item \textbf{Change support:} Helped teams adopt structured workflows, reproducible practices, versioned methods and clearer validation logic.
  \item \textbf{Software platform:} Developed CATpyLISA, a Python platform supporting workflow automation, model comparison, validation and technical review.
  \textcolor{DarkPastelRed}{\href{https://gitlab.oca.eu/gboileau/lisa_l3_tools}{\bf GitLab:CATpyLISA}}
\end{itemize}

\divider

\cvevent{\textbf{Data Scientist - Performance Analysis, Diagnostics \& Decision Support}}{Universiteit Antwerpen, Belgium}{2021 -- 2023}{Antwerp, Belgium}

\textbf{Context:} Analysis of precision-instrument performance in a demanding international environment involving complex data, uncertainty, diagnostics and technical reporting.

\textbf{Objective:} Identify performance drivers, analyse degradation factors and deliver structured evidence to support technical decisions.

\begin{itemize}
  \item \textbf{Performance analysis:} Analysed instrumental and environmental effects affecting sensitivity, stability and measurement quality.
  \item \textbf{Data-driven diagnostics:} Developed Python pipelines to identify dominant contributors, weak points and degradation patterns.
  \item \textbf{Decision support:} Produced structured technical analyses to help stakeholders understand limitations and prioritise mitigation actions.
  \item \textbf{Indicator analysis:} Compared measurements, simulations and statistical indicators to detect trends, anomalies and performance changes.
  \item \textbf{Reporting:} Delivered reproducible reports, technical summaries and presentations for international engineering and research stakeholders.
  \item \textbf{Collaboration:} Worked with multidisciplinary teams in a transnational environment requiring clear communication and rigorous documentation.
\end{itemize}

\divider

\cvevent{\textbf{Junior R\&D Engineer - Process Modelling, Simulation \& Technical Analysis}}{ARTEMIS, Observatoire de la C\^ote d'Azur}{2018 -- 2021}{Nice, France}

\textbf{Context:} Engineering-oriented research for gravitational-wave mission preparation, involving simulation workflows, large datasets, uncertainty analysis and system-performance assessment.

\textbf{Objective:} Build robust analysis methods and support structured investigations of complex technical systems.

\begin{itemize}
  \item \textbf{Complex analysis:} Developed methods for separating overlapping low-SNR signals in large and complex datasets.
  \item \textbf{Workflow development:} Built simulation and analysis workflows for uncertainty studies, performance assessment and result comparison.
  \item \textbf{Process reliability:} Compared model behaviours and outputs to assess consistency, robustness and limiting factors.
  \item \textbf{Technical investigations:} Investigated noise sources and correlations through diagnostic analyses relevant to precision instruments.
  \item \textbf{Scientific communication:} Produced peer-reviewed results and communicated complex findings to expert communities.
\end{itemize}


\cvsection{Professional Training}

\cvevent{Supervised Machine Learning (Regression \& Classification)}{DeepLearning.AI - Andrew Ng}{2020}{Coursera}

\divider

\cvevent{Recherche reproductible : principes m'ethodologiques pour une science transparente}{FUN MOOC -- Inria}{2025}{France Universit'e Num'erique}

\divider

\cvevent{Continuous Professional Training -- Software, Data, AI and Systems}{OpenClassrooms}{2024 -- 2026}{}

\small{
Project-oriented digital skills: Python, SQL, Machine Learning, AI, Linux, JavaScript, TypeScript, Angular, Node.js, Django, REST APIs, C/C++, web development, algorithms and software engineering fundamentals.
}

\switchcolumn

\cvsection{Profile for PRO BTP}

\textbf{Project-oriented engineer with strong experience in transformation of technical workflows, cross-functional coordination, process structuring, performance analysis, reporting and decision support}. Background developed in complex international environments, complemented by recent industrial software experience involving digital services, operational processes and stakeholder coordination.

For a \textbf{Chef de Projet Organisation} role, I bring a structured and analytical approach: analysing complex information, formalising needs, improving processes, defining target workflows, coordinating stakeholders, producing decision-support material and supporting change adoption.

\begin{itemize}
  \item Strong ability to analyse functional, technical and organisational information and translate it into clear recommendations.
  \item Experience structuring workflows, processes, validation logic, documentation and coordination routines across multiple stakeholders.
  \item Ability to identify operational weaknesses, inefficiencies, risks, quality issues and improvement opportunities.
  \item Experience producing structured reporting, progress summaries, technical notes and review material for decision-making.
  \item Strong analytical and synthesis skills developed through PhD-level work, international projects and complex system analysis.
  \item Comfortable working at the interface between technical contributors, project stakeholders, management needs and operational constraints.
  \item Practical experience with collaborative tools including Confluence, Zenhub, GitLab and structured documentation workflows.
  \item Strong command of \textbf{PowerPoint}, \textbf{Excel} and \textbf{Word} for executive synthesis, reporting, indicator follow-up and governance-oriented deliverables.
\end{itemize}

\cvsection{Fit for Organisation \& Transformation}

\begin{itemize}
  \item Relevant experience for transverse project coordination, process improvement, workflow transformation and performance-oriented analysis.
  \item Ability to frame problems, consolidate needs, analyse indicators and formalise decision-support material.
  \item Experience defining target workflows, improving documentation practices and supporting adoption of structured working methods.
  \item Strong fit for reporting, stakeholder alignment, technical synthesis and preparation of governance-oriented deliverables using \textbf{PowerPoint}, \textbf{Excel} and \textbf{Word}.
  \item Comfortable in high-standard environments requiring rigor, reliability, autonomy, listening skills and structured communication.
  \item Motivated to apply project coordination, analytical and transformation skills to organisational performance and process improvement challenges.
\end{itemize}

\cvsection{Languages}

\cvskill{English}{4}
\divider
\cvskill{Spanish}{2}
\divider

\medskip

\cvsection{Education}

\cvevent{PhD in Physics}{Universit\'e C\^ote d'Azur}{2018 -- 2021}{Nice, France}

\divider

\cvevent{Selective Graduate Program in Fundamental Physics (Magist\`ere)}{Universit\'e Paris-Saclay}{2016 -- 2018}{Orsay, France}

\divider

\cvevent{MSc in Astronomy and Astrophysics}{Observatoire de Paris / Universit\'e Paris-Saclay}{2017 -- 2018}{Paris, France}

\divider

\cvevent{BSc in Fundamental Physics}{Universit\'e Paris-Saclay}{2015 -- 2016}{Orsay, France}

\end{paracol}


\cvsection{Projects}

\cvevent{CNES / LISA -- Workflow Transformation, Project Coordination \& Performance Analysis}{CNES, Observatoire de la C\^ote d'Azur / Laboratoire Lagrange}{2023 -- 2025}{Nice, France}

Coordination and structuring of software, data-analysis and validation workflows for the LISA ESA/NASA space mission, within an international collaboration involving ESA, NASA and multiple research institutes.

\textbf{Project management relevance:} Structured priorities, technical tasks, validation workflows, documentation practices and contributor alignment across a multi-stakeholder environment.

\textbf{Organisation relevance:} Work involving process structuring, traceability, reporting, consistency checks, risk identification, technical synthesis, workflow improvement and decision-support material.

\textbf{Deliverables:} Python platform, validation summaries, technical notes, progress updates, documented workflows and review material supporting project execution and technical governance.

\textbf{Program scale:} Major international space mission with an estimated budget of approximately 2 billion EUR over 10 years.

\divider

\cvevent{Einstein Telescope and LIGO--Virgo--KAGRA -- Performance Studies \& Technical Coordination}{Universiteit Antwerpen}{2021 -- 2023}{Antwerp, Belgium}

Contribution to data analysis, signal characterization, performance diagnostics and technical investigations within large-scale international collaborations connected to gravitational-wave detector systems.

\textbf{Project management relevance:} Coordinated technical analyses, structured reproducible workflows, produced reporting material and communicated results to international research and engineering stakeholders.

\textbf{Organisation relevance:} Work involving indicator analysis, performance limitations, prioritisation of mitigation directions, documentation quality, technical synthesis and cross-functional communication.

\textbf{System impact:} Studies on environmental and instrumental noise sources helping identify constraints, limitations and improvement directions for next-generation observatories.

\cvsection{Strengths}

\vspace{-.3cm}

\cvsubsection{Project organisation \& transformation}

{\small
\cvtag{Project Management}
\cvtag{Transformation Projects}
\cvtag{Organisation Analysis}
\cvtag{Process Improvement}
\cvtag{Process Mapping}
\cvtag{Target Process Definition}
\cvtag{Change Management}
\cvtag{Change Support}
\cvtag{Stakeholder Alignment}
\cvtag{Governance Support}
\cvtag{Decision Support}
\cvtag{Roadmap Support}
\cvtag{Risk Identification}
\cvtag{Action Plan Follow-Up}
\cvtag{Workshop Facilitation}
}

\divider\smallskip
\vspace{-.3cm}

\cvsubsection{Analysis, reporting \& performance}

{\small
\cvtag{Functional Analysis}
\cvtag{Technical Analysis}
\cvtag{Performance Analysis}
\cvtag{Indicator Analysis}
\cvtag{Reporting}
\cvtag{Dashboard Logic}
\cvtag{Executive Synthesis}
\cvtag{Technical Synthesis}
\cvtag{Requirements Formalisation}
\cvtag{Documentation Quality}
\cvtag{Quality Improvement}
\cvtag{Operational Efficiency}
}

\divider\smallskip
\vspace{-.3cm}

\cvsubsection{Digital, SI \& tools}

{\small
\cvtag{Digital Projects}
\cvtag{Information Systems}
\cvtag{Data Workflows}
\cvtag{Software Workflows}
\cvtag{MS Office}
\cvtag{PowerPoint}
\cvtag{Excel}
\cvtag{Word}
\cvtag{Executive Presentations}
\cvtag{Indicator Follow-Up}
\cvtag{Python}
\cvtag{SQL}
\cvtag{Confluence}
\cvtag{Jira}
\cvtag{Zenhub}
\cvtag{GitLab}
\cvtag{GitHub}
\cvtag{Docker}
\cvtag{REST API}
\cvtag{Technical Documentation}
}

\divider\smallskip
\vspace{-.3cm}

\cvsubsection{Soft skills}

{\small
\cvtag{Analytical Thinking}
\cvtag{Synthesis Skills}
\cvtag{Organisation}
\cvtag{Reliability}
\cvtag{Autonomy}
\cvtag{Listening Skills}
\cvtag{Dialogue}
\cvtag{Mediation}
\cvtag{Cross-Functional Communication}
\cvtag{Adaptability}
\cvtag{Execution Rigor}
\cvtag{Technical English}
}

\end{document}
